% Chapter 3 — System Design and Architecture

\chapter{System Design and Architecture}
\label{Chapter3}

This chapter presents the high-level design and architecture of Kent-AI. It describes the multi-phase processing pipeline, the inter-component data flows, the technology stack decisions, and the design principles that govern the system's structure. The architecture is designed for modularity, testability, and extensibility.

%----------------------------------------------------------------------------------------

\section{System Overview}

Kent-AI is a multi-phase clinical NLP pipeline that transforms an unstructured conversational patient narrative into a structured, ranked, and physician-verifiable clinical report. The system comprises six major functional layers:

\begin{enumerate}
  \item \textbf{Conversational Intake Layer}: Finite-state machine guiding patient through 7-dimension clinical case-taking.
  \item \textbf{NER Layer}: Multi-encoder transformer ensemble extracting structured entity spans.
  \item \textbf{Reasoning Layer}: LLM-based clinical resolver for negation, coreference, and query synthesis.
  \item \textbf{Retrieval Layer}: Dense semantic search over 74,513 rubrics via ChromaDB HNSW.
  \item \textbf{Ranking Layer}: Grade-weighted totality scorer implementing Kent's prescribing hierarchy.
  \item \textbf{Presentation Layer}: Doctor-in-the-loop Streamlit workstation with PDF export.
\end{enumerate}

\subsection{End-to-End Data Flow}

Figure \ref{fig:architecture} illustrates the end-to-end data flow through the pipeline:

\begin{figure}[h]
\centering
\begin{minipage}{0.9\textwidth}
\begin{verbatim}
  PATIENT CONSULTATION
         |
         v
  [Intake FSM — 10 States x 7 Dimensions]
         |  Conversation Transcript
         v
  [Multi-Encoder NER — 4 Models, 15 BIO Labels]
         |  Raw Entity Spans
         v
  [LLM Clinical Resolver — Negation, Coreference]
         |  Structured 7-Dim Symptom Profile
         v
  [ChromaDB HNSW — 74,513 Rubric Index, <18ms]
         |  Top-K Rubric Candidates
         v
  [Grade-Weighted Totality Ranker — Kent's 3-Tier]
         |  Ranked Remedy Report
         v
  [Streamlit Workstation + ReportLab PDF]
\end{verbatim}
\end{minipage}
\caption[Kent-AI End-to-End Pipeline]{End-to-end data flow through the Kent-AI processing pipeline.}
\label{fig:architecture}
\end{figure}

%----------------------------------------------------------------------------------------

\section{Component Architecture}

\subsection{Conversational Intake Engine (src/chatbot/)}

The intake layer implements a \textbf{10-State Finite State Machine (FSM)} that guides the consultation through each of the seven Kentian symptom dimensions. The FSM is defined in \code{src/chatbot/state\_machine.py} with the \code{IntakeState} enumeration.

Key design decisions:
\begin{itemize}
  \item \textbf{Adaptive Slot Skipping}: If a patient's response implicitly covers multiple dimensions, the FSM advances past already-satisfied states, reducing consultation length.
  \item \textbf{Denied Slots}: Dimensions explicitly denied by the patient (``I have no headache'') are marked as denied and excluded from retrieval queries.
  \item \textbf{Deterministic Fallbacks}: If the LLM is unavailable (Ollama offline), the system degrades gracefully to rule-based template responses defined in \code{src/chatbot/fallbacks.py}.
  \item \textbf{Three-Tier Slot Extraction}: The \code{slot\_extractor.py} module implements a cascade: (1) LLM extraction, (2) pattern-based extraction, (3) conservative null fallback.
\end{itemize}

The \textbf{Dynamic Suggestion Generator} (\code{src/chatbot/suggestions.py}) streams context-aware quick-reply chips to the UI, using an LLM intent classifier to generate relevant suggestions (e.g., \textit{``It has been ongoing for a while''}, \textit{``It comes and goes in episodes''}) anchored to the current intake state.

\subsection{Synthetic Data Generation Layer (src/data/)}

The training data pipeline comprises:

\begin{itemize}
  \item \textbf{Case Generator} (\code{case\_generator.py}): LLaMA 3 8B via Ollama REST API generates structured JSON cases with narrative text and BIO-annotated entity spans.
  \item \textbf{BIO Auto-Tagger} (\code{bio\_tagger.py}): Implements exact character-offset tokenization with subword alignment repair, achieving zero offset drift across 99,854 annotated entities.
  \item \textbf{Stratified Splitter} (\code{splitter.py}): 80/10/10 train/val/test split with deficit-balancing algorithm ensuring each rubric appears in all partitions.
\end{itemize}

\subsection{Multi-Encoder NER Layer (src/models/)}

The NER layer supports four encoder architectures through a unified \code{SymptomNER} interface:

\begin{table}[h]
\caption[NER Model Architecture Summary]{Summary of the four fine-tuned biomedical encoder architectures.}
\label{tab:models}
\centering
\begin{tabular}{lllr}
\toprule
\tabhead{Model} & \tabhead{Base Checkpoint} & \tabhead{Pre-training Corpus} & \tabhead{Params} \\
\midrule
Bio\_ClinicalBERT & \texttt{emilyalsentzer/Bio\_ClinicalBERT} & MIMIC-III clinical notes & 110M \\
PubMedBERT & \texttt{microsoft/BiomedNLP-BiomedBERT} & PubMed abstracts + full text & 110M \\
BioLinkBERT & \texttt{michiyasunaga/BioLinkBERT-base} & PubMed + citation graph links & 110M \\
DeBERTa-v3 & \texttt{microsoft/deberta-v3-base} & General domain + RTD objective & 184M \\
\bottomrule
\end{tabular}
\end{table}

The \textbf{4-Model Majority Voting Ensemble} (\code{src/models/ensemble.py}) collects raw BIO predictions from all four models and resolves conflicts using weighted IoU span clustering. The \textbf{Two-Stage Hybrid Pipeline} (\code{src/models/hybrid\_pipeline.py}) uses the ensemble as a candidate miner and passes ambiguous spans to a LLaMA 3.1 LLM arbitrator for resolution.

\subsection{LLM Clinical Reasoning Layer (src/models/resolver.py)}

The \code{SymptomResolver} module invokes LLaMA 3 via the Ollama REST API in JSON mode to:
\begin{enumerate}
  \item Classify each extracted span into one of the 7 Kentian dimensions.
  \item Apply negation heuristics (``no headache'', ``denies anxiety'').
  \item Resolve coreference (``it'' referring to previously mentioned pain location).
  \item Synthesize ChromaDB query strings optimized for rubric retrieval.
\end{enumerate}

The \textbf{Structured JSON schema} enforces that every output conforms to the \code{SymptomProfile} dataclass, with rule-based fallback parsing for malformed LLM outputs.

\subsection{QLoRA Large Language Model Pipeline (src/models/)}

Two large language models were fine-tuned using QLoRA PEFT:

\begin{itemize}
  \item \textbf{BioMistral 7B} (\code{BioMistral/BioMistral-7B}): Biomedical LLM fine-tuned on 15,786 instruction-formatted training cases. 4-bit NF4 quantization, rank-16 LoRA adapters targeting all attention projection matrices. Adapter checkpoint: 81MB.
  \item \textbf{Qwen 2.5 7B} (\code{Qwen/Qwen2.5-7B-Instruct}): General instruction LLM fine-tuned identically. Adapter checkpoint: 78MB.
\end{itemize}

Both models are deployed via the \texttt{llm\_dataset.py} instruction dataset transformer which formats training examples as \code{<|system|> ... <|user|> {narrative} <|assistant|> {json\_extraction}} templates.

\subsection{Dense Retrieval Layer (src/search/)}

The retrieval layer comprises two modules:
\begin{itemize}
  \item \textbf{RubricEmbedder} (\code{embedder.py}): Lazy-loaded \texttt{all-MiniLM-L6-v2} sentence transformer producing 384-dimensional unit-normalized embeddings.
  \item \textbf{RubricVectorStore} (\code{vector\_store.py}): ChromaDB persistent HNSW cosine index over 74,513 rubrics. Supports section filtering (e.g., restrict to MIND chapter only), similarity score thresholds, and batch upsert.
\end{itemize}

The index was built in 931 seconds at 80.0 rubrics/second on the NVIDIA RTX 5000 Ada. At query time, Top-K retrieval runs in under 18 milliseconds.

\subsection{Remedy Ranking Layer (src/search/ranker.py)}

The \code{RemedyRanker} implements Kent's hierarchical prescribing formula:

\begin{equation}
\text{Score}(R) = \sum_{r \in \text{Matched Rubrics}} \text{Grade}(R, r) \times \text{Sim}(r) \times \text{Spec}(r)
\label{eq:ranking}
\end{equation}

where:
\begin{itemize}
  \item $\text{Grade}(R, r) \in \{3.0, 2.0, 1.0\}$ is Kent's bibliographic grade for remedy $R$ under rubric $r$ (Bold = 3, Italic = 2, Roman = 1).
  \item $\text{Sim}(r)$ is the cosine similarity score from ChromaDB retrieval.
  \item $\text{Spec}(r) = \log(N / \text{freq}(r))$ is the inverse rubric frequency specificity, penalizing over-broad rubrics.
\end{itemize}

Dimension weights further modulate the score: Mental symptoms (MENT) contribute 3×, Modalities (MOD) 2×, Sensations (SEN) 2×, and Location (LOC) 1×.

\subsection{Clinical Workstation (src/dashboard/)}

The Streamlit workstation implements a \textbf{4-step consultation workflow}:

\begin{enumerate}
  \item \textbf{Step 1 — Patient Intake}: Three modes: Guided Conversation (FSM chatbot), Paste Clinical Notes (manual text extraction), Enter Manually (dimension-by-dimension form). The 7:3 column layout places the chat console in 70\% width with the Case Snapshot rail in 30\%.
  \item \textbf{Step 2 — Review and Edit}: Clinician reviews and corrects extracted entities, validates dimension coverage, and approves the symptom profile.
  \item \textbf{Step 3 — Analysis}: Interactive Totality Matrix (Rubrics × Remedies, Grade badges), candidate remedy comparison, rubric Mute/Eliminatory toggles.
  \item \textbf{Step 4 — Report}: Consolidated clinical report with physician brief, Matplotlib totality score chart, ReportLab PDF export.
\end{enumerate}

Additional navigation tabs include: Case Library, Rubric Explorer (hybrid semantic/keyword search over 74,513 rubrics), Materia Medica (Grade-3 keynote browser), Methodology, and Settings.

%----------------------------------------------------------------------------------------

\section{Technology Stack}

Table \ref{tab:techstack} summarises the technology stack.

\begin{table}[h]
\caption[Technology Stack]{Kent-AI Technology Stack.}
\label{tab:techstack}
\centering
\begin{tabular}{lll}
\toprule
\tabhead{Layer} & \tabhead{Technology} & \tabhead{Version / Details} \\
\midrule
Language & Python & 3.12 \\
NER Framework & HuggingFace Transformers & 4.46+ \\
PEFT / QLoRA & HuggingFace PEFT + BitsAndBytes & 0.14+ \\
LLM Runtime & Ollama & LLaMA 3 8B / 3.1 8B \\
Vector Store & ChromaDB & 0.5+ \\
Embeddings & sentence-transformers & all-MiniLM-L6-v2 \\
Training & PyTorch & 2.6.0 (CUDA 12.4) \\
Evaluation & seqeval & Token BIO F1 \\
Dashboard & Streamlit & $\geq$1.64.0 \\
PDF Export & ReportLab + PyMuPDF & 5.1+ \\
Database & SQLite & cases.sqlite, repertory.sqlite \\
Test Framework & pytest & 195 tests, 22 modules \\
Hardware & NVIDIA RTX 5000 Ada & 32GB VRAM, CUDA 12.4 \\
\bottomrule
\end{tabular}
\end{table}

%----------------------------------------------------------------------------------------

\section{Design Principles}

The codebase adheres to the following engineering principles, documented in \code{Context/CONVENTIONS.md}:

\begin{itemize}
  \item \textbf{Modularity}: Every functional concern is encapsulated in a dedicated module under 500 lines. No monolithic files.
  \item \textbf{Type Safety}: \code{from \_\_future\_\_ import annotations} and comprehensive Python type hints on all public interfaces. \code{ClinicalCase} is the single typed source of truth for case data.
  \item \textbf{Testability}: Every module has a corresponding test file. Mock backends (\code{--mock} flag) enable fast unit testing without GPU or Ollama dependency.
  \item \textbf{Reproducibility}: Random seeds set across NumPy, PyTorch, and Python \code{random} in all training scripts.
  \item \textbf{Security}: All HTML rendered in the Streamlit frontend is sanitized via a mandatory \code{esc()} utility function.
  \item \textbf{Accessibility}: UI components use ARIA roles, Material Icon set (no emoji), and SR-only live regions for screen reader support.
\end{itemize}
